% Generated draft: requires mathematical and source review.
\documentclass[11pt]{article}
\usepackage[margin=25mm]{geometry}
\usepackage{fontspec}
\setmainfont{Latin Modern Roman}
\usepackage{amsmath,amssymb,mathtools,unicode-math}
\setmathfont{Latin Modern Math}
\usepackage{enumitem,needspace,xurl,hyperref}
\setlength{\parindent}{0pt}
\setlength{\parskip}{5pt}
\setlength{\emergencystretch}{3em}
\newcommand{\N}{\mathbb{N}}
\newcommand{\Z}{\mathbb{Z}}
\newcommand{\Q}{\mathbb{Q}}
\newcommand{\R}{\mathbb{R}}
\newcommand{\C}{\mathbb{C}}
\newcommand{\F}{\mathbb{F}}
\newcommand{\A}{\mathbb{A}}
\newcommand{\PP}{\mathbb{P}}
\newcommand{\E}{\mathbb{E}}
\newcommand{\eps}{\varepsilon}
\newcommand{\dd}{\,\mathrm{d}}
\begin{document}
\section*{UnsolvedMath ID: 3100051}
\textbf{Source problem code:} catalogue entry
\subsection{Definitions and mathematical statement}
For a real number \(x\), write \(\{x\}=x-\lfloor x\rfloor\in[0,1)\) for its fractional part, and put \(x_n=\{(3/2)^n\}\) for every positive integer \(n\). The question is whether this sequence is uniformly distributed modulo \(1\): does
\[
\lim_{N\to\infty}\frac{1}{N}\#\{1\le n\le N:a\le x_n<b\}=b-a
\qquad(0\le a<b\le1)
\]
hold for every such interval? Equivalently, if \(r_n\) is the unique integer with \(0\le r_n<2^n\) and \(r_n\equiv3^n\pmod {2^n}\), the question concerns the uniform distribution of \(r_n/2^n\).

The supplied research summary labels this problem open and says that even density of \(\{x_n:n\ge1\}\) in \([0,1]\) is unknown. It attributes infinitely many accumulation points in each of \([0,1/2]\) and \([1/2,1]\) to Pisot (1938) and Vijayaraghavan (1941), and attributes a lower bound of \(1/3\) for the diameter of the closure to Flatto--Lagarias--Pollington. It also reports a more general bound \(1/p\) for a rational base \(p/q\), without stating the hypotheses on \(p,q\). These are supplied background claims, not resolutions of the uniform-distribution question. The summary mentions Beukers--Dubickas estimates for separation from integers, with the number \(0.5769\), but supplies no precise inequality to formalize; it also cautions that the cited Furstenberg-type dynamics are different. No unspecified estimate or generalization is adopted as part of the statement.

The source page's open label belongs to its October 2020 snapshot. Its later resolution was not checked item by item in the supplied extraction, so the registry's status is retained as dated source information.
\subsection{Short English statement}
Do the fractional parts of \((3/2)^n\), for positive integers \(n\), visit every interval with limiting frequency equal to that interval's length? The supplied record labels this open and reports that even density is unknown; these status statements and the partial results are source attributions.
\subsection{Sources}
\begin{itemize}
\item Local registry: \url{data/unsolvedmath/problems.json}, permanent UnsolvedMath ID \(3100051\), source code catalogue entry; complete supplied statement, background, and research summary, updated 2026-08-06.
\item Cooper, \emph{Combinatorial Problems I Like}, October 2020 snapshot, Combinatorial Number Theory, source-order bullet 51: \url{https://people.math.sc.edu/cooper/combprob.html}. HTML extraction accessed 2026-07-29. The original statement's ``See this'' link has no target in the supplied record.
\item Background attributions supplied by the registry: Pisot (1938); Vijayaraghavan (1941); Flatto--Lagarias--Pollington; Beukers--Dubickas; and Furstenberg-type ergodic methods. No full bibliographic references or exact theorem statements for these attributions were supplied.
\end{itemize}
\end{document}
